\section{Conclusion}
\label{sec:conclusion}

This paper verified, through a three-stage experiment on four lightweight VLMs
and three public medical VQA datasets, the full process of adapting 2--3B-class
vision--language models to the medical imaging VQA domain on consumer-grade GPUs:
zero-shot baselines (Phase~1), 75 QLoRA fine-tuning conditions (Phase~2), and an
810-trial programme comparing hyperparameter search strategies and re-testing
configuration consistency (Phase~3).

For RQ1 the null hypothesis was rejected: the correctness patterns of the four
models differed significantly both on a pooled basis ($n = 8{,}231$; Cochran's
$Q = 1904.28$, $df = 3$, $p < .001$) and on all three datasets individually.
Qwen3-VL-2B achieved the highest pooled accuracy at 0.3843, though close enough
to Qwen2.5-VL-3B to be statistically indistinguishable on some datasets, while
Gemma4-E2B was significantly below all three other models; the ranking was
preserved after removing samples suspected of pretraining exposure. For RQ2,
whether the null hypothesis was rejected depended on the model: performance
improved significantly for Qwen2.5-VL-3B ($d = +2.646$) and Qwen3-VL-2B
($d = +1.620$), whereas SmolVLM2-2.2B significantly deteriorated ($d = -2.284$)
and Gemma4-E2B was not significant ($d = -0.652$), and the pooled mixed-effects
estimate of no significant effect ($p = .3629$) reflects the cancellation of
opposing model-level effects rather than the absence of an effect. For RQ3 the
null hypothesis was rejected in the direction opposite to the one it anticipated:
Optuna (0.4490) was significantly superior to Autoresearch (0.4184)
($U = 16.00$, $p = .0112$, $r = -0.68$), and Autoresearch was statistically
indistinguishable even from Random Search (0.4186), its lower-bound reference,
while all three automated strategies exceeded the manual configuration (0.3776).
The second requirement of RQ3, interpretable search rationales, was not tested by
the four-strategy comparison, because the system prompt instructed the agent to
respond with JSON only; 147 of 200 trials (73.5\%) returned only the
hyperparameter JSON. In the corrected re-experiment all 200 responses carried a
natural-language rationale (100\%, mean length 620 characters), so that rationales
can be produced at all is confirmed, but no qualitative evaluation of their
interpretability was performed. Search quality was meanwhile poor: 12.5 unique
configurations out of 20 trials on average, recovering to 18.90 / 20 once the
configuration inconsistencies were removed, which eliminated the significant
deficit against Optuna ($p = .0112 \rightarrow .1726$) --- the run-level mean
nevertheless still trailed Optuna (0.4292 vs.\ 0.4490) and improvement over the
original was not established ($p = .2387$), so the negative answer to the first
requirement of RQ3 stands.

The contributions are three. First, three-stage empirical data on medical-domain
adaptation of lightweight VLMs under a single evaluation protocol, including the
finding that performance and resource consumption are not monotonically related
and the indirect result that on SLAKE closed-ended accuracy a QLoRA-adapted 2B
model can reach parity with a 7B medical-specialized model (85.26 vs.\ 85.34)
while a gap of 8--11\,\%p remains on PathVQA and VQA-RAD. Second, a
quantification of the model-level heterogeneity of fine-tuning effects, showing
that an identical procedure can simultaneously produce large improvement and
significant deterioration depending on the model and that estimating these
together leads to the misreading of ``no effect.'' Third, a negative result for
LLM-based autonomous HPO reported together with a verification of its failure
mechanism and an experimental settlement of the question of correctability: the
duplicate proposals were empirically shown to be a product of prompt--
implementation configuration inconsistency rather than an intrinsic limitation of
the agent, which indicates that future work on LLM-based HPO must treat
prompt--implementation consistency as a controlled variable. Alongside these, we
record two cases in which the choice of aggregation unit reversed the conclusion,
and the quantitative counterexample that the variation observed when repeating an
identical configuration (about 0.056) exceeds the mean difference between
strategies detected here (0.031).

Four directions follow. Autonomous HPO should be re-verified with an expanded
search budget, since for Autoresearch --- in the original condition and in v2
alike --- the upper bound of the IQR for the trial of best performance touches the
budget ceiling of 20; and the consistency-restored condition in particular needs
both a larger sample (an effect of magnitude $r = +0.32$ is not detectable at
$n = 10$) and a search policy that reliably converts secured diversity into
high-performance attainment, together with a decision criterion distinguishing
stochastic variation across re-runs from genuine improvement. A qualitative
evaluation framework for search rationales is needed, measuring causal
correspondence between a rationale and the actual proposal, post-hoc
rationalization, and expert acceptability. Improving open-ended response
performance is the most urgent task, and building an evaluation set that
deliberately oversamples the clinically important types is a prerequisite for
judging it. Finally, a re-experiment controlling effective training volume, and
comparison with medical-specialized VLMs under an identical protocol together
with a clinically validated evaluation scheme, remain to be done.
